\documentclass[11pt]{article}
\usepackage[letterpaper,left=.82in,right=.82in,top=.72in,bottom=.7in]{geometry}
\usepackage[T1]{fontenc}
\usepackage{lmodern,microtype,amsmath,amssymb,booktabs,tabularx,array}
\usepackage{xcolor,listings,enumitem,fancyhdr,titlesec}
\usepackage[hidelinks]{hyperref}
\hypersetup{pdftitle={jwselect: Quick User Guide, 0.5.0},pdfauthor={Doosoo Kim}}
\definecolor{navy}{HTML}{233D54}
\definecolor{pale}{HTML}{F3F5F7}
\newcommand{\cmd}[1]{\texttt{#1}}
\newcolumntype{Y}{>{\raggedright\arraybackslash}X}
\setlength{\parindent}{0pt}
\setlength{\parskip}{4pt}
\setlength{\emergencystretch}{2em}
\renewcommand{\arraystretch}{1.10}
\setlist[itemize]{leftmargin=1.3em,itemsep=2pt,topsep=3pt}
\titleformat{\section}{\large\bfseries\color{navy}}{\thesection}{.6em}{}
\titlespacing*{\section}{0pt}{10pt}{5pt}
\titleformat{\subsection}{\normalsize\bfseries}{\thesubsection}{.5em}{}
\titlespacing*{\subsection}{0pt}{7pt}{3pt}
\lstset{basicstyle=\ttfamily\small,columns=fullflexible,keepspaces=true,
 breaklines=true,showstringspaces=false,backgroundcolor=\color{pale},
 frame=none,xleftmargin=6pt,xrightmargin=6pt,aboveskip=5pt,belowskip=5pt}
\pagestyle{fancy}
\fancyhf{}
\fancyhead[L]{\small\color{navy}\texttt{jwselect} \textbar\ Quick user guide}
\fancyhead[R]{\small Version 0.5.0}
\fancyfoot[R]{\small\thepage}
\renewcommand{\headrulewidth}{.3pt}
\setlength{\headheight}{14pt}
\begin{document}
\thispagestyle{empty}
{\LARGE\bfseries\color{navy}\texttt{jwselect}: Quick User Guide}\par
{\normalsize Doosoo Kim (\href{mailto:doosoonet@gmail.com}{doosoonet@gmail.com})}\par
{\small Version 0.5.0 \hfill 26 October 2026}\par
\vspace{3pt}
Select optional preperiod restrictions and estimate treatment effects by
minimum distance (MD), using cohort--time estimates and their joint covariance.
This guide covers the ordinary estimation workflow, including optional CS-DR
and nonlinear DiD adapters in version 0.5.0.

\section{Install}
The main command requires Stata 16.1 or later. Extract the current archive and
replace the path below with the folder containing \cmd{jwselect.pkg}:
\begin{lstlisting}
net install jwselect, from("C:/path/to/jwselect") replace
which jwselect
findfile _jwselect_mata.do
help jwselect
\end{lstlisting}
After an update, save your work and restart Stata. To use \cmd{jwselect}
directly after \cmd{jwdid}, first install its dependencies:
\begin{lstlisting}
ssc install require, replace
ssc install ftools, replace
ssc install reghdfe, replace
ssc install hdfe, replace
net install jwdid, ///
    from("https://raw.githubusercontent.com/friosavila/stpackages/main") replace
\end{lstlisting}
The matrix interface requires neither upstream package. The CS-DR adapter
requires \cmd{csdid} 2.x, version 2.0.1 or later. The README gives the pinned
2.0.1 install command; consult the
\href{https://psantanna.com/csdid/}{authors' installation guide} when updating
an existing installation. \cmd{jwselect} does not install dependencies.
Run commands containing \cmd{///} together from the Do-file Editor.
An optional installation check, in a fresh Stata 17 session, is
\cmd{do run\_release\_tests.do} from the extracted package folder.
The test clears data and estimates; save your work first.

\section{A typical analysis}
Replace the variable names below with yours. Code never-treated units as
\cmd{first\_treat=0}; treated units carry their first treatment date.
\begin{lstlisting}
jwdid y x1 x2, ivar(id) tvar(year) gvar(first_treat) ///
    never method(regress) cluster(id)
estimates store JW_SOURCE

* Maintained-baseline estimate, with no optional restrictions
jwselect, select(manual) postwindow(0 2)
estimates store JW_BASE

* Default: cohort-wise Wald-BIC, followed by joint MD
estimates restore JW_SOURCE
jwselect, select(wbic) candidates(nested) postwindow(0 2)
estimates store JW_WBIC
\end{lstlisting}
\textbf{Restore \cmd{JW\_SOURCE} before every new direct \cmd{jwselect} run.}
Each successful run replaces Stata's estimation results. A bare \cmd{jwselect}
only redisplays the current results. Here, \cmd{postwindow(0 2)} averages the
first three treatment periods within each cohort.

\textbf{Supported linear input.} Use a balanced, regularly spaced panel,
unweighted linear estimation, full cohort--time effects, clustering by
\cmd{id}, and numeric covariates constant within each unit. The default
\cmd{vce(unconditional)} includes covariate-averaging uncertainty and requires
\cmd{method(regress)}. For other designs, use the matrix interface on page~\pageref{sec:matrix}.

\section{CS-DR and nonlinear DiD}
\subsection*{CS-DR: start from unaggregated \texttt{csdid} results}
\begin{lstlisting}
csdid y x1 x2, ivar(id) time(year) gvar(first_treat) ///
    method(dr) nevertreated base_period(universal) ///
    analytical pointwise cluster(id)
estimates store CS_SOURCE
jwselect, select(wbic) postwindow(0 2)
estat selection
estimates restore CS_SOURCE
jwselect, select(manual) postwindow(0 2)
\end{lstlisting}
Use an unweighted balanced panel with positive integer, unit-spaced time,
never-treated units, and plain numeric covariates constant within unit.
Clustering can be omitted or specified at \cmd{ivar()}; higher-level clusters
are unsupported. Retain the original estimation data. Older \cmd{csdid} 1.x,
aggregated results, bootstrap covariance, varying bases, and not-yet-treated
controls are rejected. \cmd{vce(supplied)} is the default and only covariance
convention for this source; it retains the full analytical joint covariance.

\subsection*{Nonlinear DiD: logit or Poisson through \texttt{jwdid}}
\begin{lstlisting}
jwdid y x1 x2, ivar(id) tvar(year) gvar(first_treat) ///
    never method(logit) cluster(id)
estimates store NL_SOURCE
jwselect, select(wbic) postwindow(0 2)
estimates restore NL_SOURCE
jwselect, select(manual) postwindow(0 2)
\end{lstlisting}
Use \cmd{method(poisson)} for an exponential conditional mean. These two
methods use the same balanced, unweighted, never-treated, unit-clustered,
time-invariant-covariate contract as the linear adapter. The extracted
estimates are outcome-scale ATT and placebo margins, with their joint
covariance. \cmd{vce(unconditional)} is the default; \cmd{vce(conditional)}
conditions on the observed covariate distribution. Other nonlinear methods,
including \cmd{ppmlhdfe}, are outside this initial support.

\subsection*{Keep the source baseline and interpret the adjustment}
Both new adapters require \cmd{baseline(cohort)}. Omitted
\cmd{anticipation()} is inferred from the source; an explicit mismatch is
rejected. For \cmd{jwdid}, the compatible value is its recorded anticipation
minus one. Restore the stored source before every new direct specification.

Optional restrictions set \emph{average outcome-scale placebo effects} to
zero. The estimator applies a Gaussian covariance-weighted moment adjustment;
it does not refit the nonlinear model or re-estimate CS-DR with alternative
bases. The Wald--$J$ identity refers to this quadratic MD problem. With
covariate interactions, zero mean placebo effects need not set every index
pretrend coefficient to zero. Validity requires the relevant identifying,
nuisance-model and joint-covariance conditions for every moment used.
Adjusted nonlinear effects are not clipped to natural outcome bounds.
Fixed \cmd{jwselect} aggregation weights need not match \cmd{csdid}'s defaults.

The self-contained files \path{examples/csdid_demo.do} and
\path{examples/nonlinear_demo.do} illustrate both workflows.

\newpage
\section{Choose the design and preperiod restrictions}
\subsection*{Baseline and treatment-effect horizon}
The baseline is maintained, never selected or pretrend-tested. Optional dates add restrictions.
Common-baseline choices below apply to linear and matrix inputs; the two new
adapters retain their upstream cohort baseline.
For annual data, $g$ is the cohort's treatment date and $q$ the earliest input treatment date.
With \cmd{baseline(common)} and no anticipation, eligible dates extend through $g-1$,
excluding the baseline $q-1$.

\begin{tabularx}{\linewidth}{@{}p{.38\linewidth}Y@{}}
\toprule
Option & Meaning \\
\midrule
\cmd{baseline(cohort)} & Baseline $g-1$; the default. \\
\cmd{baseline(common)} & Baseline $q-1$ for every cohort. \\
\cmd{anticipation(1)} & Allow one anticipation period: use $g-2$ or $q-2$ and exclude potentially affected optional dates. Default: 0 for linear/matrix inputs; inferred for the new adapters. \\
\cmd{postwindow(0 0)} & Immediate effect only, for each cohort. \\
\cmd{postwindow(0 2)} & Average over event times 0, 1, and 2. \\
\bottomrule
\end{tabularx}

Without \cmd{postwindow()}, all available post-treatment cells are used.
Every requested date must exist for every cohort. For spacing $h=\cmd{gap()}$
and anticipation $\delta$, the baseline is $g-h(1+\delta)$ or $q-h(1+\delta)$;
event time is $(t-g)/h$. Direct input obtains spacing from the source; the
CS-DR adapter currently requires unit spacing.
The two commands use different anticipation conventions; do not copy the
\cmd{jwdid} anticipation value mechanically into \cmd{jwselect}.

\subsection*{Baseline only, manual choice, or automatic selection}
\begin{tabularx}{\linewidth}{@{}p{.38\linewidth}Y@{}}
\toprule
Option & What it does \\
\midrule
\cmd{select(manual)} & Use only the maintained baseline unless \cmd{prematrix()} is supplied. \\
\cmd{select(wbic)} & Default. Select within each cohort; add \cmd{wbicscope(joint)} for joint selection (next page). \\
\cmd{select(focused)} & Balance estimated bias and variance for a chosen target; see page~\pageref{sec:focused}. \\
\bottomrule
\end{tabularx}

For manual selection, list optional dates in a two-column matrix:
\emph{cohort, calendar time}. For example, if cohorts 4 and 6 and these dates
are present in your data:
\begin{lstlisting}
matrix P = (4,1 \ 4,2 \ 6,3 \ 6,4)
estimates restore JW_SOURCE
jwselect, select(manual) prematrix(P) postwindow(0 2)
\end{lstlisting}
Do not list the baseline. Dates need not be consecutive. Manual and Wald--BIC
use the \emph{union} of the chosen restrictions in joint MD: one cohort's
restrictions can affect other estimates through covariance. A cohort absent
from \cmd{P} contributes no optional restrictions of its own.

\subsection*{Choose the candidate family explicitly}
With optional dates $\{1,2,3\}$, \cmd{candidates(nested)} searches
$\varnothing,\{3\},\{2,3\},\{1,2,3\}$.
\cmd{candidates(all)} searches all eight subsets, including $\{1,3\}$.
With $K$ optional dates, these families have $K+1$ and $2^K$ candidates.
The package default is \cmd{nested}; request \cmd{all} explicitly for exhaustive search.
Baseline only is included in both families.

The limits are \cmd{maxcand(4096)} per cohort and \cmd{maxjoint(4096)} for
a joint WBIC or aggregate-focused family. Exceeding a limit causes an error;
no candidates are silently dropped.
All-subset searches grow quickly; raise a limit only if you intend the larger family.

\textbf{Aggregation.} The default gives cohorts equal weight, divided equally
among their included post periods. \cmd{aggregate(equalcell)} weights included
cells equally.

\subsection*{Joint Wald--BIC selection (new in 0.4.0)}
Use one joint criterion to choose restrictions across all cohorts:
\begin{lstlisting}
estimates restore JW_SOURCE
jwselect, select(wbic) wbicscope(joint) ///
    candidates(nested) postwindow(0 2)
estat selection
estat scores
estimates store JW_JOINT
\end{lstlisting}
The joint criterion is $W_A-r_A\log N$: the joint pretrend Wald statistic,
minus the number of restrictions times the log of the independent-unit
count. It uses covariance across cohorts and dates, then estimates all
effects by joint MD. The default \cmd{wbicscope(cohort)} selects within
cohorts before joint MD; these two rules can choose different sets.

The joint family combines every local candidate with every candidate in
other cohorts. With $K_g$ optional dates in cohort $g$, there are
$\prod_g(K_g+1)$ nested configurations or $2^{\sum_g K_g}$ all-subset
configurations. Even nested joint search may exceed \cmd{maxjoint(4096)}.
A prespecified \cmd{jointmatrix(JM)} can supply a smaller family: one binary
joint mask per row, with columns in \cmd{e(hmap)} order. Obtain that map with
\cmd{estat design}; see \cmd{help jwselect} for details. Fix the family before
inspecting outcomes. Raise a cap only for a computationally feasible family.

\subsection*{Joint pretrend tests}
After any successful \cmd{jwselect} fit:
\begin{lstlisting}
* All eligible optional preperiod restrictions, jointly
estat pretrend
return list

* Manual/Wald-BIC selected set: diagnostic after selection
estat pretrend, selected
\end{lstlisting}
For focused results, explicitly choose an output instead:
\begin{lstlisting}
matrix list e(output_map)
estat pretrend, selected output(1)
\end{lstlisting}
\cmd{output()} means a result-row number and is allowed only with
\cmd{selected}. Manual/Wald--BIC selected tests default to row 1, since all
rows share the same joint restriction set. For a prespecified subset, use
\cmd{estat pretrend, prematrix(P)}; \cmd{P} has the same two-column
cohort/date format as manual selection. Do not combine it with \cmd{selected}.

The test reports Wald, the equivalent MD $J$, degrees of freedom, and a
chi-square reference p-value. $J$ is evaluated through the Wald identity
under the compatible covariance-weighted MD model, not by an independent
refit. \cmd{r(table)} has columns \cmd{df Wald J pvalue}; separate scalars
include \cmd{r(chi2)}, \cmd{r(J)}, \cmd{r(df)}, and \cmd{r(p)}.

\enlargethispage{\baselineskip}
\textbf{Interpretation.} A prespecified test uses the chi-square reference
under the maintained model. A p-value for an outcome-selected set is a
diagnostic, not an automatic correction for selection. Rejection concerns
optional restrictions; the maintained baseline is never tested.
For a focused aggregate of cohort-selected effects, the test uses the union
of their restrictions; $J$ refers to that set's compatible joint MD model.
An empty set has zero degrees of freedom and no p-value.

\newpage
\section{Focused estimates and confidence intervals}
\label{sec:focused}
Request cohort-focused estimates, their aggregate, and numerical inference:
\begin{lstlisting}
estimates restore JW_SOURCE
jwselect, select(focused) focus(cohort) ///
    candidates(all) postwindow(0 2) level(95)
estat inference
\end{lstlisting}
\begin{tabularx}{\linewidth}{@{}p{.29\linewidth}Y@{}}
\toprule
Option & Selection target and reported estimates \\
\midrule
\cmd{focus(cohort)} & Select for each cohort's mean effect; also report the aggregate of the resulting cell estimates. \\
\cmd{focus(aggregate)} & Select a joint restriction configuration for the aggregate. Cohort results use that same selected model. \\
\cmd{focus(both)} & Package default. Report cohort-focused effects, their aggregate, and a separately aggregate-focused estimate. \\
\bottomrule
\end{tabularx}
For point estimates only, add \cmd{inference(none)}; intervals and p-values
will be missing.

\subsection*{How to read focused inference}
The default \cmd{inference(auto)} repeats selection inside joint Gaussian
draws and searches over plausible preperiod violations. Use the reported
intervals and p-values. \cmd{se\_fixed} is a fixed-selection diagnostic;
do not build ordinary normal intervals from it. Intervals are marginal for
each effect, not simultaneous bands, and need not be narrower than baseline-only intervals.

\subsection*{Search budget: the controls that matter}
\begin{tabularx}{\linewidth}{@{}p{.29\linewidth}p{.12\linewidth}Y@{}}
\toprule
Option & Default & Meaning \\
\midrule
\cmd{optmaxeval()} & 400 & Evaluations per objective, shared across starts. \\
\cmd{optstarts()} & 4 & Starting points per objective. \\
\cmd{optreps()} & 1,999 & Common Gaussian draws used during search. \\
\cmd{reps()} & 19,999 & Independent draws used to check discovered points. \\
\cmd{opttol()} & 0.02 & Stopping mesh size in standardized coordinates. \\
\cmd{seed()} & 12,345 & Reproducible simulation seed. \\
\bottomrule
\end{tabularx}

There are four objectives per effect: two interval bounds and two p-value
tails. Each start receives 100 evaluations by default.
\textbf{``Search budget exhausted'' means at least one start reached its allowance.}
The command returns the best points found; this is not a Stata failure.
Inspect \cmd{budget\_exhausted} and \cmd{converged\_starts} in \cmd{estat inference}.

First increase the evaluation budget while retaining the same design:
\begin{lstlisting}
estimates restore JW_SOURCE
jwselect, select(focused) focus(cohort) ///
    candidates(all) postwindow(0 2) level(95) optmaxeval(1600)
estat inference
\end{lstlisting}
Compare intervals and p-values with the earlier run. More starts, more draws,
and another seed provide further checks. Increasing \cmd{reps()} alone does
not resolve budget exhaustion. Converged searches do not certify global extrema.

The defaults \cmd{gamma(.01)} and \cmd{mcerr(.001)} allocate error to the
nuisance region and simulation checks. See the full guide before changing them.

\newpage
\section{Read, save, and supply results}
\subsection*{Inspect and save a completed run}
Each cohort row averages its included post periods. The aggregate uses the
chosen cell weights. \cmd{e(output\_map)} identifies the actual cohort behind
each row; an aggregate has cohort value zero.
\begin{lstlisting}
jwselect                    // Redisplay the current results
estat selection             // Candidate sets and selected restrictions
estat design                // Baselines and cohort-time maps
estat cells                 // Individual cohort-time estimates
matrix list e(output_map)
estimates save "jwselect_results.ster", replace
\end{lstlisting}
These commands follow a successful \cmd{jwselect} call. Store or save each
result you want to compare before restoring \cmd{JW\_SOURCE}.

\begin{tabularx}{\linewidth}{@{}p{.38\linewidth}Y@{}}
\toprule
Stored result & Use \\
\midrule
\cmd{e(table)} & Main table: estimate, \cmd{se\_fixed}, interval bounds, and zero-effect p-value. \\
\cmd{e(b)}, \cmd{e(ci)}, \cmd{e(pvalue)} & Estimates, confidence intervals, and p-values separately. \\
\cmd{e(selection)} & Restrictions used by each output; absent when no optional preperiods exist. \\
\cmd{e(V)} & Joint covariance for manual and Wald--BIC results. \\
\cmd{e(V\_fixed)} & Fixed-selection covariance; diagnostic for focused estimates. \\
\bottomrule
\end{tabularx}

Source metadata record the estimator, version, method, model, and restriction
scale. Inspect them with \cmd{ereturn list}. The source version is empty when
unavailable, including \cmd{jwdid}; retain its \cmd{which} output separately.

Prespecify manual restrictions, candidate families, and target weights.
Manual intervals treat the chosen restrictions as fixed. Wald--BIC intervals
use a pointwise oracle approximation under the method's assumptions; they
are not uniformly valid near violations. Focused results do not post an
ordinary \cmd{e(V)}. Do not use \cmd{e(V\_fixed)} with conventional
\cmd{lincom} or \cmd{test} as if it adjusted for focused selection.

\subsection*{Alternative input: coefficient and covariance matrices}
\label{sec:matrix}
Use this route for separately validated cohort--time estimates and their full
covariance, including designs outside the direct adapters. It requires neither
\cmd{jwdid} nor \cmd{csdid}.
\begin{lstlisting}
jwselect, bmat(B) vmat(V) cells(C) nunits(500) ///
    gap(1) select(wbic) candidates(all) postwindow(0 2)
\end{lstlisting}
\begin{tabularx}{\linewidth}{@{}p{.23\linewidth}Y@{}}
\toprule
Input & Required contents \\
\midrule
\cmd{B} & Row or column vector of unpooled treated-cohort--time estimates. \\
\cmd{V} & Full sampling covariance of \cmd{B}, including covariance across cohorts and dates. Supply $V$, not $NV$. \\
\cmd{C} & Matching rows: \emph{cohort, calendar time}. Include the normalization cells and only positive treated-cohort codes. \\
\cmd{nunits()} & Number of independent sampled units, used in the Wald--BIC penalty; replace 500 with your count. \\
\cmd{gap()} & Period spacing, with default 1. \\
\bottomrule
\end{tabularx}

\enlargethispage{\baselineskip}
Repeat the matrix arguments for each specification. The convention is
\cmd{vce(supplied)}. Supply unpooled cohort--time estimates and a validated full
covariance; retain structural-zero normalizations, but never substitute zeros
for unidentified effects.

\end{document}
